% XeLaTeX chapter fragment; see review.tex for a minimal compile entry point.
\chapter{公开知识图谱问答实验与结果分析}
\label{chap:public-experiments}

\section{实验问题与评价协议}

本章围绕第\ref{chap:execution-feedback}章的方法回答三个问题：恢复轨迹训练是否优于同监督预算的普通续训；这种差异能否在另一续训随机种子及未用于调参的问题上保留；额外的强化学习、详细报错反馈与推理成本分别提供何种证据。主要结论基于最后一次冻结的项目留出评测，开发集结果用于说明方法选择和扩展试验的边界。

\subsection{数据来源与划分边界}

实验使用 KQA Pro 及其知识库执行环境。项目从官方训练数据中划分 $5\,000$ 道训练题、$500$ 道开发题和 $2\,000$ 道项目留出题。划分按完全相同问题或完全相同程序形成的连通分组进行，排除项目历史已使用的分组及与官方验证集完全重复的问题或程序分组。冻结记录显示，新划分之间完全相同问题和程序的重叠数均为零。

此划分降低了项目内部的直接重复使用问题，但有明确边界。项目留出集\textbf{不是官方隐藏测试集}，也不是按全部组合结构严格隔离的组合泛化测试；排除完全重复不能消除所有近似题、共享实体和共享知识。基础模型预训练是否见过数据同样无法排除。此前已接触的官方验证数据只作为历史参考，不用于本章最终留出结论。

在访问项目留出结果前，固定五个模型、原始反馈模式、解码配置、主统计量、成本口径与描述性分组。五路生成全部结束后，先在不读取答案的条件下回放四组智能体轨迹及重新执行程序基线；共 $10\,000$ 条记录全部通过，零差异。随后打开冻结答案文件完成统一评分。该回放检查保证记录与执行一致，不保证自然语言条件完整或答案语义正确。评测结束后不再基于这批留出题调参或选择最佳随机种子。

\subsection{对照模型}

\begin{description}
  \item[P：一次性程序生成。] 由相同基础模型生成完整查询程序，再统一执行。使用通过参考回放核验的训练问题，训练三个遍历。它提供不同问答范式的效果与成本参照。
  \item[I：初始分步模型。] 基于成功的分步执行轨迹监督训练两个遍历，用于采集训练自由运行轨迹及初始化后续训练。
  \item[A：普通配对续训。] 从 I 出发，以普通成功轨迹和不插入真实分歧动作的参考后缀续训。
  \item[C：恢复轨迹续训。] 从同一个 I 出发，在匹配问题、记录及监督 token 预算下，用含真实分歧状态的配对后缀替换 A 的后缀部分。
  \item[B、D：强化学习扩展。] 分别从第一续训种子的 A、C 继续执行相同规模的短程 GRPO 试验，只在开发集分析其增量作用。
\end{description}

A/C 是检验本方法的主要受控对照。P 的输出表示、训练遍历和推理流程均不同，P 与 C 的差异不能全部归因为恢复训练。I 不用于替代 A 作主要方法对照，因为 C 还包含额外续训。两组续训种子以下简写为下标 $1$（$20261003$）和 $2$（$20261004$）。留出集只评价预先固定的 P、$A_1$、$C_1$、$A_2$、$C_2$，没有按照开发集最高准确率选择单个模型。

\subsection{指标与不确定性}

每题以固定答案比较函数计算正确性 $y_{i,s}^{M}\in\{0,1\}$，其中 $M$ 为模型组，$s$ 为续训种子。主统计量为先在同一问题上平均两个种子的 C/A 正确性差异，再对问题平均：
\begin{align}
  d_i &= \tfrac12\big[(y_{i,1}^{C}-y_{i,1}^{A})
                    +(y_{i,2}^{C}-y_{i,2}^{A})\big],\\
  \widehat\Delta &= \frac{100}{N}\sum_{i=1}^{N}d_i,
  \qquad N=2\,000.
\end{align}
置信区间采用问题配对 bootstrap：以问题为单位有放回抽样 $5\,000$ 次，每次同时保留该问题在两种子下的所有结果，计算百分位区间。统计单位为 $2\,000$ 个问题，不能把两个种子当成 $4\,000$ 个独立观测。区间条件于已训练出的检查点，未覆盖从基础模型开始的全流程训练随机性，也不校正开发过程中反复查看结果带来的选择影响。

除准确率外，记录完成率、调用次数、无效调用、输入 token、生成 token 及二者总和。总 token 按每轮未填充的完整输入加生成 token 累计，不能直接等同于浮点计算量、实际延迟或费用。对照中的“改正”表示基线错误而候选正确，“退步”表示基线正确而候选错误；二者差值对应净增正确题数。执行失败或未完成的题目不会从分母中删除。

\section{项目留出集主要结果}

\begin{table}[htbp]
  \centering
  \caption{$2\,000$ 题项目留出集结果。每个模型使用完整的相同题集。}
  \label{tab:holdout-main}
  \PublicHoldoutTable
\end{table}

\begin{table}[htbp]
  \centering
  \caption{项目留出集配对差异。区间单位为百分点，条件于对应检查点。}
  \label{tab:holdout-paired}
  \PublicPairedTable
\end{table}

两个种子的 A 平均准确率为 $78.75\%$，C 为 $83.50\%$，主统计量为提升 $4.75$ 个百分点，$95\%$ 配对区间为 $[3.53,5.98]$。$C_1$ 相对 $A_1$ 提高 $4.65$ 个百分点，改正 $177$ 题、退步 $84$ 题；$C_2$ 相对 $A_2$ 提高 $4.85$ 个百分点，改正 $171$ 题、退步 $74$ 题。两次续训均得到正向净收益，且各自区间均高于零。

这些结果支持以下限定结论：在本项目的基础模型、执行环境、训练数据和推理预算下，恢复状态与后缀配对训练的整体方案在匹配监督 token 的普通续训对照上具有稳定的留出准确率优势。它没有证明所有问题均改善；两次仍分别存在 $84$ 和 $74$ 道由正确转为错误的题。两种子共享初始模型，因此结果也不能写成“经过两次完全独立训练验证”。

P 的准确率为 $76.15\%$。两个 C 分别比 P 高 $6.75$ 和 $7.95$ 个百分点，但由于训练表示和推理交互方式不同，这属于方法范式参照，不能作为恢复模块的单独因果效应。

\section{执行稳定性与成本}

\begin{table}[htbp]
  \centering
  \caption{留出集累计推理成本；无效调用是调用级计数，未完成是题级计数。}
  \label{tab:holdout-cost}
  \resizebox{\textwidth}{!}{\PublicCostTable}
\end{table}

$A_1$ 和 $A_2$ 分别有 $198$、$205$ 题未完成，C 降至 $80$、$91$ 题；无效调用分别减少约 $65.96\%$ 和 $68.37\%$。这与训练方案改善执行稳定性的解释一致，但它仍是整体训练差异后的观测，不能单凭这些指标判定唯一作用机制。

相对对应的 A，两个 C 的总 token 分别减少约 $8.84\%$ 和 $4.17\%$，主要体现累计输入开销下降。生成 token 却分别增加约 $3.88\%$ 和 $16.35\%$。因此不应概括成所有资源消耗都下降，也不能在没有等硬件延迟分析的情况下声称普遍加速。

与一次性程序基线 P 相比，$C_1$ 和 $C_2$ 的总 token 分别为 P 的约 $23.64$ 和 $24.72$ 倍。分步推理需要反复处理历史上下文，这是提高准确率时必须呈现的成本取舍。本文的主要实验证据是相对于匹配分步基线的质量改进，尚不支持优于一次性程序生成的全面效率结论。

\section{开发集复现与扩展试验}

\subsection{续训随机种子的配对复现}

\begin{table}[htbp]
  \centering
  \caption{$500$ 题开发集的 A/C 配对结果。}
  \label{tab:dev-pair}
  \PublicDevTable
\end{table}

第一续训种子的开发集收益为 $2.40$ 个百分点，区间包含零；第二续训种子的收益为 $5.20$ 个百分点，区间高于零。只依据第一轮开发结果不足以形成稳定的效果结论。因此保留相同方案，追加配对续训，并在协议冻结后统一评价两个种子的留出结果。开发集在方案推进中已被查看，其结果作为开发记录，不能与留出结果合并成额外独立验证。

\subsection{短程强化学习没有给出明确增量收益}

第一续训种子的 A/C 各开展一次同规模 GRPO 扩展，得到 B/D。每组 $200$ 次更新，使用相同的 $200$ 道训练题，每题四条采样轨迹，共 $800$ 条；开发评价保持原有严格推理预算。训练日志确认存在有效梯度和参数变化，故该试验并非只有流程运行而没有训练更新。

\begin{table}[htbp]
  \centering
  \caption{开发集强化学习扩展的配对结果。B/D 分别由 $A_1$/$C_1$ 续训得到。}
  \label{tab:rl-extension}
  \PublicRLTable
\end{table}

B 的准确率为 $78.00\%$，相对 A 的 $78.40\%$ 下降 $0.40$ 个百分点；D 为 $82.00\%$，相对 C 的 $80.80\%$ 提高 $1.20$ 个百分点，两个增量区间均包含零。D 相对 C 的总 token 还增加约 $2.38\%$。两项 RL 增量均未达到此前设定的继续投入标准：准确率提高至少 $2$ 个百分点，或准确率损失不超过 $1$ 个百分点且总 token 至少减少 $20\%$。

D 相对 B 的 $4.00$ 个百分点差异是\textbf{经过 RL 后仍保留的恢复训练组差异}，不能写成 RL 带来的提升。本项目据此停止扩大 RL 试验，保留 C 为主要方案。这一结果只说明该单种子、短程配置没有提供明确的额外收益，不说明强化学习在此类任务上必然无效。

\subsection{详细报错文字的机制证据不足}

为检验模型是否利用详细执行错误说明，在固定开发题和固定 A/C 检查点上进行反馈呈现干预：仅把失败观察中的错误类型和细节替换为通用失败文本，保留成功结果、句柄、成功标志及真实执行状态。该诊断既不重新训练，也不把屏蔽版本作为留出模型候选。

\begin{table}[htbp]
  \centering
  \caption{详细报错文字屏蔽诊断；每组 $500$ 题，差值为屏蔽减原始。}
  \label{tab:feedback-diagnostic}
  \PublicFeedbackTable
\end{table}

四组区间均包含零，未观察到稳定的详细报错文字优势。全部 $4\,000$ 条原始和屏蔽轨迹通过执行回放，但 C 中出现五条在该题自身首次干预前就发生输出差异的轨迹。有限批次复核显示，同题此前的提示和事件一致，而同批其他题的干预会改变批次组成或填充长度；数值批处理差异是一种可能解释，因缺少受控重跑与 logits 证据，尚未证实其原因。故该诊断不宜承担精确的机制因果结论。

恢复训练配对中仅 $4/418$ 的分歧动作被拒绝，这项干预覆盖的是执行反馈中的狭窄部分。结论应为“详细报错文字的作用没有得到证实”，不能据此否定整个恢复状态训练，也不能声称已经证明模型会理解语义错误。最终留出评测继续使用原始反馈。

\section{预设问题组的描述性分析}

\begin{table}[htbp]
  \centering
  \caption{预设问题组的 C/A 准确率差异。分组可以重叠，仅作描述。}
  \label{tab:descriptive-groups}
  \PublicGroupTable
\end{table}

不同长度与操作类型的分组中，两次续训的 C/A 点估计均为正。限定条件组共有 $483$ 题，差异较大，分别为 $11.59$ 和 $8.70$ 个百分点。这个现象可作为雷达领域条件查询研究的动机，但不能直接表述为“条件忠实性显著提升”：这里按参考程序类型分组，评价仍然是最终答案比较，未逐条人工审查生成路径是否保留问题全部条件。各组存在重叠，也未进行独立子组显著性或多重比较推断。

\section{有效性限制与论文结论范围}

\begin{enumerate}
  \item \textbf{数据独立性有限。} 留出题只相对于本项目已知的训练、开发和历史使用做了去重隔离，不能排除基础模型的预训练接触，也不代表严格组合泛化。
  \item \textbf{随机性覆盖有限。} 两种子共享初始适配器、配对语料和开发题。问题 bootstrap 只估计条件于检查点的题目抽样不确定性。
  \item \textbf{受控变量仍是整体训练方案。} A/C 匹配监督 token，但输入长度和计算量不相同；尚未分离额外状态、句柄变化与后缀监督的独立作用。
  \item \textbf{答案正确不等于语义忠实。} 执行回放和标准答案比较能够检查不同层面的正确性，但无法独立证明执行路径遵守了全部限定条件。
  \item \textbf{训练恢复状态覆盖有限。} $418$ 个配对来自有限的失败轨迹，重复使用不能增加独立状态种类。公开知识库和雷达领域的模式、值类型、语言及来源要求仍有差异。
  \item \textbf{成本优势有条件。} C 相对 A 的累计总 token 减少，但相对 P 仍约为 $24$ 倍；未给出统一成本约束下优于所有基线的结论。
  \item \textbf{基线范围有限。} 本章使用可运行且来源一致的项目基线，未复现全部已有公开方法。结果支持本训练方案相对该基线的改进，不能作为公开排行榜领先声明。
\end{enumerate}

在上述范围内，公开实验证实恢复轨迹配对监督值得保留为论文主要方法。它的可见收益包括更高的最终问答准确率、更少的未完成题目和无效调用；详细报错文字及短程强化学习的额外价值尚无充分证据。下一阶段重点转向雷达领域的事实核验、可执行查询与独立评价，而非继续利用已打开的项目留出集搜索更高分数。

\section{结果可追溯性与后续材料}

本章表格由冻结的公开汇总报告生成，导出程序只读取汇总 JSON 与源码摘要，不读取逐题预测、留出答案或模型权重。\texttt{results\_tables.csv} 给出数值、原始字段与派生公式，\texttt{evidence\_manifest.json} 绑定来源文件及章节的 SHA-256。准确率、配对净增和成本比例可由汇总计数复算；已冻结的 bootstrap 区间按来源复制并校验文件摘要，跨种子联合区间不能仅凭边际计数重新计算。

本章未引入尚未核对的文献条目。正式学位论文仍需结合学校模板统一章节编号，补充经核实的 KQA Pro、基础模型、LoRA、分步问答与恢复训练相关文献，并完成雷达独立评价章节。当前两章是新主线下已具备实验依据的写作材料，不表示全文或雷达实验已经完成。
